\documentclass[../main]{subfiles}
\ifSubfilesClassLoaded{\setcounter{chapter}{11}}{}
\begin{document}

\chapter{Metode Pembuktian Langsung dan Kontraposisi}

\begin{subcpmk}
  \item \textbf{Sub-CPMK 4.1:} Mahasiswa mampu menerapkan berbagai metode pembuktian matematis dasar (pembuktian langsung, pembuktian tidak langsung/kontradiksi, dan kontraposisi).
\end{subcpmk}

% ============================================================
% MATERI POKOK
% ============================================================
\section{Pengertian Teorema, Bukti, dan Terminologi}

Dalam matematika dan ilmu komputer, setiap klaim mengenai sifat suatu objek, algoritma, atau sistem harus divalidasi secara formal melalui proses pembuktian. Pembuktian bukan sekadar verifikasi empiris (menguji beberapa kasus), melainkan demonstrasi logis yang menunjukkan bahwa suatu pernyataan berlaku untuk \emph{seluruh} kasus yang mungkin. Kemampuan menyusun dan memahami bukti merupakan kompetensi fundamental dalam ilmu komputer, karena pembuktian menjadi dasar bagi analisis algoritma, verifikasi program, dan kriptografi \cite{rosen2019}.

\subsection{Teorema}

\begin{definisi}[Teorema]
\logic{Teorema} adalah pernyataan matematika yang kebenarannya telah dibuktikan secara formal melalui rangkaian penalaran logis yang valid, berdasarkan aksioma, definisi, dan teorema-teorema lain yang telah dibuktikan sebelumnya. Teorema merupakan kebenaran yang berlaku secara universal dalam sistem aksiomatik yang mendasarinya.
\end{definisi}

Sebagian besar teorema memiliki struktur implikasi: ``Jika $p$, maka $q$'' ($p \rightarrow q$), di mana $p$ disebut \logic{hipotesis} (premis, anteseden) dan $q$ disebut \logic{kesimpulan} (konklusi, konsekuen). Membuktikan teorema berarti menunjukkan bahwa implikasi $p \rightarrow q$ bernilai benar untuk setiap kasus --- yaitu tidak pernah terjadi situasi di mana $p$ benar tetapi $q$ salah.

\subsection{Bukti (Proof)}

\begin{definisi}[Bukti]
\logic{Bukti} (\textit{proof}) adalah demonstrasi deduktif yang memaparkan langkah demi langkah penalaran logis --- berdasarkan aksioma, definisi, dan aturan inferensi yang sah --- untuk menunjukkan bahwa suatu proposisi atau teorema bernilai benar. Setiap langkah dalam bukti harus dijustifikasi oleh landasan logis yang valid \cite{velleman2019}.
\end{definisi}

Bukti berbeda secara fundamental dari verifikasi empiris. Seorang ilmuwan mungkin menguji hipotesis dengan mengamati ribuan kasus dan menemukan bahwa semuanya sesuai, tetapi hal itu tidak membuktikan bahwa hipotesis berlaku untuk \emph{seluruh} kasus. Sebaliknya, bukti matematis memberikan jaminan absolut bahwa suatu pernyataan benar untuk setiap elemen dalam domain yang dimaksud.

\subsection{Terminologi Terkait}

Beberapa konsep penting yang erat kaitannya dengan teorema dan bukti:

\begin{enumerate}[label=\textbf{\arabic*.}]
  \item \textbf{Aksioma / Postulat}: Pernyataan kebenaran mendasar yang diterima tanpa pembuktian. Aksioma menjadi titik awal dari mana teorema-teorema diturunkan. Contoh: ``Untuk setiap bilangan bulat $a$, $a + 0 = a$'' (identitas penjumlahan).

  \item \textbf{Definisi}: Pernyataan yang secara tepat menentukan makna suatu istilah atau konsep matematis. Definisi bukan klaim yang perlu dibuktikan, melainkan konvensi yang disepakati. Contoh: ``Bilangan bulat $n$ disebut genap jika terdapat bilangan bulat $k$ sehingga $n = 2k$''.

  \item \textbf{Lemma}: Teorema bantu yang dibuktikan terlebih dahulu dan digunakan sebagai batu loncatan untuk membuktikan teorema utama yang lebih besar. Lemma memecah pembuktian yang kompleks menjadi bagian-bagian yang lebih mudah dikelola.

  \item \textbf{Akibat (\textit{Corollary})}: Proposisi yang kebenarannya mengikuti secara langsung dan mudah dari teorema yang telah dibuktikan, tanpa memerlukan pembuktian tambahan yang substansial.

  \item \textbf{Konjektur (\textit{Conjecture})}: Pernyataan yang diyakini benar berdasarkan bukti empiris atau intuisi, tetapi belum dibuktikan secara formal. Contoh terkenal: Konjektur Goldbach (setiap bilangan genap $> 2$ adalah jumlah dua bilangan prima).
\end{enumerate}

\begin{konsep}
Hierarki dalam matematika dapat digambarkan sebagai: \textbf{Aksioma} $\rightarrow$ \textbf{Lemma} $\rightarrow$ \textbf{Teorema} $\rightarrow$ \textbf{Akibat}. Aksioma diterima tanpa bukti, lemma dibuktikan untuk mendukung teorema, teorema dibuktikan sebagai hasil utama, dan akibat diturunkan dari teorema dengan mudah.
\end{konsep}

\subsection{Notasi Q.E.D.}

Akhir dari sebuah bukti biasanya ditandai dengan simbol $\blacksquare$ atau singkatan \textbf{Q.E.D.} (\textit{quod erat demonstrandum}, bahasa Latin untuk ``yang hendak dibuktikan''). Simbol ini mengindikasikan bahwa rangkaian penalaran telah lengkap dan kesimpulan yang dikehendaki telah tercapai. Dalam konteks buku ini, simbol $\blacksquare$ akan digunakan secara konsisten untuk menandai akhir setiap bukti.

\section{Pembuktian Langsung (Direct Proof)}

Pembuktian langsung merupakan teknik pembuktian yang paling fundamental dan intuitif. Untuk membuktikan implikasi $p \rightarrow q$, kita mengasumsikan bahwa hipotesis $p$ bernilai benar, kemudian melalui serangkaian langkah penalaran deduktif yang valid, menunjukkan bahwa kesimpulan $q$ juga bernilai benar. Teknik ini memanfaatkan fakta bahwa implikasi $p \rightarrow q$ bernilai salah hanya jika $p$ benar dan $q$ salah; dengan menunjukkan bahwa $q$ selalu benar ketika $p$ benar, kita membuktikan bahwa kasus palsu tersebut tidak mungkin terjadi \cite{rosen2019}.

\subsection{Struktur Bukti Langsung}

Prosedur umum pembuktian langsung untuk implikasi $p \rightarrow q$:

\begin{enumerate}
    \item \textbf{Asumsikan} hipotesis $p$ bernilai benar.
    \item \textbf{Terapkan} definisi, aksioma, lemma, dan teorema yang berlaku.
    \item \textbf{Lakukan} manipulasi aljabar dan penalaran logis selangkah demi selangkah.
    \item \textbf{Tunjukkan} bahwa rangkaian penalaran tersebut menghasilkan kesimpulan $q$ bernilai benar.
\end{enumerate}

Setiap langkah dalam bukti harus dijustifikasi secara eksplisit. Justifikasi dapat berupa: penerapan definisi, penggunaan aksioma, hasil dari lemma atau teorema sebelumnya, atau aturan inferensi yang valid.

\subsection{Contoh: Kuadrat Bilangan Ganjil}

\begin{contoh}
\textbf{Teorema}: Jika $n$ bilangan ganjil, maka $n^2$ adalah bilangan ganjil.

\textbf{Bukti (Langsung):}
\begin{enumerate}
    \item Asumsikan $n$ adalah bilangan ganjil. \hfill (\textit{asumsi hipotesis})
    \item Berdasarkan definisi bilangan ganjil, terdapat bilangan bulat $k$ sehingga $n = 2k + 1$. \hfill (\textit{definisi})
    \item Hitung $n^2$:
    \[ n^2 = (2k+1)^2 = 4k^2 + 4k + 1 = 2(2k^2 + 2k) + 1 \]
    \item Misalkan $m = 2k^2 + 2k$. Karena $k$ bilangan bulat, maka $m$ juga bilangan bulat. \hfill (\textit{ketertutupan})
    \item Sehingga $n^2 = 2m + 1$, yang berbentuk definisi bilangan ganjil.
    \item Kesimpulan: $n^2$ adalah bilangan ganjil. $\blacksquare$
\end{enumerate}
\end{contoh}

\subsection{Contoh: Jumlah Dua Bilangan Genap}

\begin{contoh}
\textbf{Teorema}: Jika $a$ dan $b$ bilangan genap, maka $a + b$ juga bilangan genap.

\textbf{Bukti (Langsung):}
\begin{enumerate}
    \item Asumsikan $a$ dan $b$ bilangan genap. \hfill (\textit{asumsi})
    \item Berdasarkan definisi, terdapat bilangan bulat $j$ dan $k$ sehingga $a = 2j$ dan $b = 2k$. \hfill (\textit{definisi genap})
    \item Hitung $a + b$:
    \[ a + b = 2j + 2k = 2(j + k) \]
    \item Misalkan $m = j + k$, yang merupakan bilangan bulat (karena penjumlahan dua bilangan bulat menghasilkan bilangan bulat).
    \item Sehingga $a + b = 2m$, yang berbentuk definisi bilangan genap.
    \item Kesimpulan: $a + b$ bilangan genap. $\blacksquare$
\end{enumerate}
\end{contoh}

\subsection{Contoh: Hasil Kali Bilangan Rasional}

\begin{contoh}
\textbf{Teorema}: Jika $r$ dan $s$ bilangan rasional, maka $r \cdot s$ juga bilangan rasional.

\textbf{Bukti (Langsung):}
\begin{enumerate}
    \item Asumsikan $r$ dan $s$ bilangan rasional.
    \item Berdasarkan definisi bilangan rasional, terdapat bilangan bulat $a, b, c, d$ dengan $b \neq 0$ dan $d \neq 0$, sehingga $r = \frac{a}{b}$ dan $s = \frac{c}{d}$.
    \item Hitung $r \cdot s$:
    \[ r \cdot s = \frac{a}{b} \cdot \frac{c}{d} = \frac{ac}{bd} \]
    \item Karena $a, c$ bilangan bulat, maka $ac$ bilangan bulat (ketertutupan perkalian). Karena $b \neq 0$ dan $d \neq 0$, maka $bd \neq 0$.
    \item Sehingga $r \cdot s = \frac{ac}{bd}$ merupakan rasio dua bilangan bulat dengan penyebut tak nol, yang memenuhi definisi bilangan rasional.
    \item Kesimpulan: $r \cdot s$ bilangan rasional. $\blacksquare$
\end{enumerate}
\end{contoh}

\begin{catatan}
Kunci keberhasilan bukti langsung terletak pada penguasaan definisi. Perhatikan bahwa pada ketiga contoh di atas, langkah pertama setelah asumsi selalu menerapkan \textbf{definisi} dari konsep yang dimaksud (ganjil, genap, rasional). Definisi memberi kita ekspresi aljabar yang dapat dimanipulasi menuju kesimpulan.
\end{catatan}

\section{Pembuktian dengan Kontraposisi}

Tidak jarang pembuktian langsung menghadapi hambatan teknis --- misalnya ketika manipulasi aljabar dari hipotesis $p$ menuju kesimpulan $q$ menjadi rumit atau tidak alami. Dalam situasi tersebut, pembuktian dengan \logic{kontraposisi} menjadi alternatif yang sangat efektif. Metode ini dilandasi oleh fakta bahwa implikasi $p \rightarrow q$ secara logis ekuivalen dengan kontraposisinya $\neg q \rightarrow \neg p$ \cite{rosen2019,forallx,libretexts_proofs}.

\subsection{Landasan Logis}

Ekuivalensi $p \rightarrow q \equiv \neg q \rightarrow \neg p$ dapat diverifikasi melalui tabel kebenaran:

\begin{center}
\renewcommand{\arraystretch}{1.2}
\begin{tabular}{|c|c||c|c|c|c||c|}
\hline
$p$ & $q$ & $\neg p$ & $\neg q$ & $p \rightarrow q$ & $\neg q \rightarrow \neg p$ & Sama? \\ \hline
T & T & F & F & T & T & Ya \\ \hline
T & F & F & T & F & F & Ya \\ \hline
F & T & T & F & T & T & Ya \\ \hline
F & F & T & T & T & T & Ya \\ \hline
\end{tabular}
\end{center}

Karena kedua kolom selalu menghasilkan nilai yang identik, keduanya merupakan ekspresi yang ekuivalen secara logis. Membuktikan kontraposisi $\neg q \rightarrow \neg p$ secara otomatis membuktikan implikasi asli $p \rightarrow q$.

\subsection{Struktur Bukti Kontraposisi}

Untuk membuktikan $p \rightarrow q$ melalui kontraposisi:

\begin{enumerate}
    \item \textbf{Asumsikan} $\neg q$ (negasi dari kesimpulan) bernilai benar.
    \item \textbf{Gunakan} penalaran deduktif berdasarkan definisi, aksioma, dan teorema yang berlaku.
    \item \textbf{Tunjukkan} bahwa $\neg p$ (negasi dari hipotesis) bernilai benar.
\end{enumerate}

\subsection{Kapan Menggunakan Kontraposisi?}

Kontraposisi sangat berguna dalam situasi-situasi berikut:
\begin{itemize}
    \item Kesimpulan $q$ mengandung negasi atau kondisi ``tidak'' yang sulit diekspresikan secara positif.
    \item Hipotesis $p$ bersifat kompleks sehingga sulit untuk memulai manipulasi aljabar darinya.
    \item Negasi dari kesimpulan ($\neg q$) memberikan informasi yang lebih mudah diolah dibandingkan hipotesis aslinya.
\end{itemize}

\subsection{Contoh: Paritas \texorpdfstring{$3n+2$}{3n+2} dan \texorpdfstring{$n$}{n}}

\begin{contoh}
\textbf{Teorema}: Jika $3n+2$ adalah ganjil, maka $n$ adalah ganjil.

\textbf{Mengapa bukti langsung sulit?} Dari $3n+2 = 2k+1$ (ganjil), didapat $n = \frac{2k-1}{3}$. Pembagian oleh 3 menyulitkan pembuktian bahwa hasilnya ganjil.

\textbf{Bukti (Kontraposisi):} Buktikan kontraposisi: ``Jika $n$ genap, maka $3n+2$ genap.''

\begin{enumerate}
    \item Asumsikan $n$ genap ($\neg q$). Maka $n = 2k$ untuk suatu bilangan bulat $k$.
    \item Hitung $3n + 2$:
    \[ 3n + 2 = 3(2k) + 2 = 6k + 2 = 2(3k + 1) \]
    \item Misalkan $m = 3k + 1$. Karena $k$ bilangan bulat, maka $m$ bilangan bulat.
    \item Sehingga $3n + 2 = 2m$, yang berbentuk bilangan genap ($\neg p$).
    \item Kontraposisi terbukti, maka teorema asli juga terbukti. $\blacksquare$
\end{enumerate}
\end{contoh}

\subsection{Contoh: Kuadrat Genap Implies Bilangan Genap}

\begin{contoh}
\textbf{Teorema}: Jika $n^2$ genap, maka $n$ genap.

\textbf{Bukti (Kontraposisi):} Buktikan: ``Jika $n$ ganjil, maka $n^2$ ganjil.''

\begin{enumerate}
    \item Asumsikan $n$ ganjil. Maka $n = 2k + 1$ untuk suatu bilangan bulat $k$.
    \item Hitung $n^2$:
    \[ n^2 = (2k + 1)^2 = 4k^2 + 4k + 1 = 2(2k^2 + 2k) + 1 \]
    \item Misalkan $m = 2k^2 + 2k$, yang merupakan bilangan bulat.
    \item Sehingga $n^2 = 2m + 1$, yang berbentuk bilangan ganjil.
    \item Kontraposisi terbukti: $n$ ganjil $\Rightarrow$ $n^2$ ganjil. Maka kontraposisinya, ``$n^2$ genap $\Rightarrow$ $n$ genap,'' juga terbukti. $\blacksquare$
\end{enumerate}
\end{contoh}

\begin{catatan}
Perhatikan bahwa contoh kedua (``$n^2$ genap $\Rightarrow$ $n$ genap'') sangat sulit dibuktikan secara langsung karena dari $n^2 = 2k$ kita memperoleh $n = \sqrt{2k}$, dan tidak ada cara langsung untuk menyimpulkan bahwa $\sqrt{2k}$ genap. Kontraposisi menghindari akar kuadrat dan membuat bukti mengalir dengan natural.
\end{catatan}

\section{Pembuktian dengan Kontradiksi (Proof by Contradiction)}

Pembuktian dengan kontradiksi (\textit{proof by contradiction}, atau \textit{reductio ad absurdum}) merupakan teknik pembuktian yang sangat kuat dan serbaguna. Metode ini bekerja dengan mengasumsikan bahwa pernyataan yang hendak dibuktikan adalah \emph{salah}, kemudian menunjukkan bahwa asumsi tersebut menghasilkan suatu kontradiksi logis --- yakni pernyataan yang mustahil benar. Karena asumsi kepalsuan menghasilkan kemustahilan, maka pernyataan asli harus benar \cite{hammack2018}.

\subsection{Landasan Logis}

Untuk membuktikan implikasi $p \rightarrow q$, bukti kontradiksi didasarkan pada ekuivalensi:
\[ p \rightarrow q \equiv \neg(p \wedge \neg q) \]

Artinya, $p \rightarrow q$ bernilai benar jika dan hanya jika $p \wedge \neg q$ bernilai salah (mustahil). Jika kita mengasumsikan $p \wedge \neg q$ (hipotesis benar, kesimpulan salah) dan berhasil menurunkan kontradiksi, maka $p \wedge \neg q$ memang mustahil, sehingga $p \rightarrow q$ terbukti benar.

Untuk membuktikan pernyataan umum $P$ (bukan implikasi), kita mengasumsikan $\neg P$ dan menurunkan kontradiksi, sehingga $\neg P$ mustahil dan $P$ harus benar.

\subsection{Struktur Bukti Kontradiksi}

\begin{enumerate}
    \item \textbf{Asumsikan} negasi dari pernyataan yang hendak dibuktikan.
    \begin{itemize}
        \item Untuk implikasi $p \rightarrow q$: asumsikan $p$ benar dan $\neg q$ benar.
        \item Untuk pernyataan umum $P$: asumsikan $\neg P$ benar.
    \end{itemize}
    \item \textbf{Lakukan} penalaran deduktif dari asumsi tersebut.
    \item \textbf{Tunjukkan} bahwa penalaran tersebut menghasilkan pernyataan yang jelas merupakan kontradiksi (misalnya: $0 = 1$, suatu bilangan genap sekaligus ganjil, atau bertentangan dengan fakta yang telah diketahui).
    \item \textbf{Simpulkan} bahwa asumsi awal (negasi) harus salah, sehingga pernyataan asli terbukti benar.
\end{enumerate}

\subsection{Contoh: Irasionalitas \texorpdfstring{$\sqrt{2}$}{sqrt(2)}}

\begin{contoh}
\textbf{Teorema}: $\sqrt{2}$ adalah bilangan irasional.

Ini merupakan salah satu bukti kontradiksi paling terkenal dalam sejarah matematika, yang telah dikenal sejak zaman Yunani kuno.

\textbf{Bukti (Kontradiksi):}
\begin{enumerate}
    \item Asumsikan, untuk mendapatkan kontradiksi, bahwa $\sqrt{2}$ adalah bilangan rasional.
    \item Maka terdapat bilangan bulat $a$ dan $b$ dengan $b \neq 0$ dan $\gcd(a, b) = 1$ (pecahan sudah disederhanakan) sehingga $\sqrt{2} = \frac{a}{b}$.
    \item Kuadratkan kedua ruas: $2 = \frac{a^2}{b^2}$, sehingga $a^2 = 2b^2$.
    \item $a^2$ genap (karena $a^2 = 2b^2$ adalah kelipatan 2), maka $a$ genap (telah dibuktikan pada section sebelumnya via kontraposisi).
    \item Karena $a$ genap, terdapat bilangan bulat $c$ sehingga $a = 2c$.
    \item Substitusi: $(2c)^2 = 2b^2$, sehingga $4c^2 = 2b^2$, maka $b^2 = 2c^2$.
    \item $b^2$ genap, maka $b$ genap.
    \item Tetapi jika $a$ dan $b$ keduanya genap, maka $\gcd(a, b) \geq 2$, yang \textbf{berkontradiksi} dengan asumsi bahwa $\gcd(a, b) = 1$.
    \item Karena kontradiksi diperoleh, maka asumsi awal salah. $\sqrt{2}$ adalah bilangan irasional. $\blacksquare$
\end{enumerate}
\end{contoh}

\subsection{Contoh: Ketidakberhinggaan Bilangan Prima}

\begin{contoh}
\textbf{Teorema} (Euclid): Terdapat tak berhingga banyak bilangan prima.

\textbf{Bukti (Kontradiksi):}
\begin{enumerate}
    \item Asumsikan, untuk mendapatkan kontradiksi, bahwa hanya terdapat berhingga bilangan prima, yaitu $p_1, p_2, \ldots, p_n$.
    \item Konstruksi bilangan $N = p_1 \cdot p_2 \cdot \ldots \cdot p_n + 1$.
    \item $N > 1$, sehingga $N$ memiliki faktor prima (setiap bilangan $> 1$ memiliki faktor prima).
    \item Misalkan $p_i$ adalah salah satu faktor prima dari $N$. Maka $p_i | N$ (membagi habis $N$).
    \item Tetapi $p_i$ juga membagi $p_1 \cdot p_2 \cdot \ldots \cdot p_n$ (karena $p_i$ adalah salah satu faktornya).
    \item Jika $p_i | N$ dan $p_i | (p_1 \cdot \ldots \cdot p_n)$, maka $p_i | (N - p_1 \cdot \ldots \cdot p_n) = 1$.
    \item Tetapi tidak ada bilangan prima yang membagi 1. \textbf{Kontradiksi.}
    \item Asumsi salah, maka terdapat tak berhingga bilangan prima. $\blacksquare$
\end{enumerate}
\end{contoh}

\subsection{Perbedaan dengan Kontraposisi}

Meskipun kontraposisi dan kontradiksi sama-sama merupakan metode tak langsung, keduanya berbeda secara fundamental:

\begin{itemize}
    \item \textbf{Kontraposisi} membuktikan pernyataan yang \emph{ekuivalen} ($\neg q \rightarrow \neg p$) secara langsung. Ini hanya berlaku untuk implikasi $p \rightarrow q$.
    \item \textbf{Kontradiksi} bekerja dengan mengasumsikan negasi dan mencari kemustahilan. Ini dapat digunakan untuk \emph{semua jenis pernyataan}, tidak terbatas pada implikasi.
\end{itemize}

\begin{konsep}
Kontradiksi sering menjadi satu-satunya pilihan untuk membuktikan pernyataan eksistensial negatif (``tidak ada ...'') atau pernyataan tentang ketidakmungkinan. Bukti bahwa $\sqrt{2}$ irasional dan bukti ketidakberhinggaan bilangan prima adalah contoh klasik di mana kontraposisi tidak dapat diterapkan secara langsung, tetapi kontradiksi memberikan pendekatan yang elegan \cite{wiki_mathematical_proof}.
\end{konsep}

\section[Pembuktian dengan Kasus]{\raggedright Pembuktian dengan Kasus dan Pembuktian Exhaustif}

Terdapat situasi di mana membuktikan suatu pernyataan secara umum (untuk seluruh elemen domain) lebih mudah dilakukan dengan membagi domain menjadi beberapa kasus yang saling lepas (\textit{mutually exclusive}) dan menangkap seluruh kemungkinan (\textit{collectively exhaustive}), kemudian membuktikan pernyataan tersebut untuk setiap kasus secara terpisah \cite{rosen2019}.

\subsection{Pembuktian dengan Kasus (Proof by Cases)}

\begin{definisi}[Pembuktian dengan Kasus]
\logic{Pembuktian dengan kasus} (\textit{proof by cases} / \textit{proof by exhaustion of cases}) adalah teknik pembuktian di mana domain dipartisi menjadi kasus-kasus $C_1, C_2, \ldots, C_n$ yang saling lepas dan lengkap, kemudian pernyataan dibuktikan secara terpisah untuk setiap kasus. Teknik ini valid berdasarkan aturan inferensi:
\[ \frac{(C_1 \rightarrow q) \wedge (C_2 \rightarrow q) \wedge \ldots \wedge (C_n \rightarrow q)}{(C_1 \vee C_2 \vee \ldots \vee C_n) \rightarrow q} \]
\end{definisi}

Kunci keberhasilan metode ini adalah memastikan bahwa kasus-kasus yang dipilih benar-benar \emph{mencakup seluruh kemungkinan}. Jika ada kasus yang terlewat, bukti menjadi tidak valid.

\begin{contoh}
\textbf{Teorema}: Untuk setiap bilangan bulat $n$, $n^2 + n$ adalah bilangan genap.

\textbf{Bukti (dengan Kasus):}

Setiap bilangan bulat $n$ bernilai genap atau ganjil (dua kasus, saling lepas dan lengkap).

\textbf{Kasus 1:} $n$ genap. Maka $n = 2k$ untuk suatu bilangan bulat $k$.
\[ n^2 + n = (2k)^2 + 2k = 4k^2 + 2k = 2(2k^2 + k) \]
yang berbentuk $2m$ (genap). $\checkmark$

\textbf{Kasus 2:} $n$ ganjil. Maka $n = 2k + 1$ untuk suatu bilangan bulat $k$.
\[ n^2 + n = (2k+1)^2 + (2k+1) = 4k^2 + 4k + 1 + 2k + 1 = 4k^2 + 6k + 2 = 2(2k^2 + 3k + 1) \]
yang berbentuk $2m$ (genap). $\checkmark$

Kedua kasus terbukti. $\blacksquare$
\end{contoh}

\subsection{Pembuktian Exhaustif}

\begin{definisi}[Pembuktian Exhaustif]
\logic{Pembuktian exhaustif} (\textit{exhaustive proof} / \textit{brute-force proof}) adalah pembuktian di mana pernyataan diverifikasi untuk \textbf{setiap} elemen domain secara individual, tanpa menggunakan argumentasi umum. Metode ini hanya dapat diterapkan jika domain bersifat berhingga dan cukup kecil untuk diperiksa satu per satu.
\end{definisi}

\begin{contoh}
\textbf{Teorema}: Untuk setiap bilangan bulat $n$ dengan $1 \leq n \leq 4$, berlaku $n^2 \leq 2^n \cdot n$.

\textbf{Bukti (Exhaustif):} Periksa setiap nilai:
\begin{center}
\renewcommand{\arraystretch}{1.2}
\begin{tabular}{|c|c|c|c|}
\hline
$n$ & $n^2$ & $2^n \cdot n$ & $n^2 \leq 2^n \cdot n$? \\ \hline
1 & 1 & 2 & Ya ($1 \leq 2$) \\ \hline
2 & 4 & 8 & Ya ($4 \leq 8$) \\ \hline
3 & 9 & 24 & Ya ($9 \leq 24$) \\ \hline
4 & 16 & 64 & Ya ($16 \leq 64$) \\ \hline
\end{tabular}
\end{center}

Semua kasus terverifikasi. $\blacksquare$
\end{contoh}

\begin{catatan}
Pembuktian exhaustif memiliki keterbatasan yang jelas: jika domain berhingga tetapi sangat besar (misalnya semua bilangan prima kurang dari $10^9$), verifikasi manual menjadi tidak praktis. Dalam kasus seperti itu, komputer sering digunakan untuk membantu verifikasi, tetapi bukti semacam itu kadang diperdebatkan statusnya dalam komunitas matematika.
\end{catatan}

\section{Penyangkalan dengan Kontra-Contoh}

Sejauh ini, pembahasan difokuskan pada cara \emph{membuktikan} bahwa suatu pernyataan benar. Namun, dalam matematika dan ilmu komputer, kemampuan untuk menunjukkan bahwa suatu pernyataan \emph{salah} sama pentingnya. Untuk menyangkal pernyataan universal (``untuk semua $x$, $P(x)$''), cukup menemukan satu elemen domain yang tidak memenuhi $P$ --- elemen tersebut disebut \logic{kontra-contoh} (\textit{counterexample}) \cite{rosen2019}.

\subsection{Landasan Logis}

Pernyataan universal $\forall x\, P(x)$ menyatakan bahwa $P(x)$ bernilai benar untuk setiap $x$ dalam domain. Negasi dari pernyataan ini adalah:
\[ \neg [\forall x\, P(x)] \equiv \exists x\, \neg P(x) \]

Artinya, untuk menyangkal pernyataan universal, cukup tunjukkan \emph{keberadaan} satu elemen $x_0$ dalam domain sehingga $P(x_0)$ bernilai salah. Elemen $x_0$ inilah yang disebut kontra-contoh.

\begin{definisi}[Kontra-Contoh]
\logic{Kontra-contoh} (\textit{counterexample}) adalah elemen spesifik dari domain yang menunjukkan bahwa suatu pernyataan universal adalah salah. Satu kontra-contoh sudah cukup untuk menyangkal pernyataan universal, tidak peduli berapa banyak contoh lain yang mendukungnya.
\end{definisi}

\subsection{Contoh Penyangkalan}

\begin{contoh}
\textbf{Klaim}: ``Untuk setiap bilangan bulat positif $n$, $n^2 - n + 41$ adalah bilangan prima.''

Klaim ini tampak benar jika diuji untuk nilai-nilai kecil:
\begin{center}
\renewcommand{\arraystretch}{1.2}
\begin{tabular}{|c|c|c|}
\hline
$n$ & $n^2 - n + 41$ & Prima? \\ \hline
1 & 41 & Ya \\ \hline
2 & 43 & Ya \\ \hline
3 & 47 & Ya \\ \hline
5 & 61 & Ya \\ \hline
10 & 131 & Ya \\ \hline
\end{tabular}
\end{center}

Namun, untuk $n = 41$:
\[ 41^2 - 41 + 41 = 41^2 = 1681 = 41 \times 41 \]
yang \textbf{bukan} bilangan prima. Kontra-contoh ditemukan pada $n = 41$, sehingga klaim tersebut \textbf{salah}. $\blacksquare$
\end{contoh}

\begin{contoh}
\textbf{Klaim}: ``Untuk setiap bilangan real $x$, $x^2 > x$.''

\textbf{Kontra-contoh}: Pilih $x = \frac{1}{2}$.
\[ x^2 = \left(\frac{1}{2}\right)^2 = \frac{1}{4} < \frac{1}{2} = x \]

Karena $x^2 < x$ (bukan $x^2 > x$), maka klaim tersebut \textbf{salah}. $\blacksquare$
\end{contoh}

\subsection{Peringatan: Kontra-Contoh vs Bukti}

\begin{konsep}
\textbf{Asimetri fundamental}: Satu kontra-contoh cukup untuk menyangkal pernyataan universal, tetapi seribu contoh yang mendukung \emph{tidak cukup} untuk membuktikan pernyataan universal. Membuktikan $\forall x\, P(x)$ memerlukan argumentasi yang berlaku untuk seluruh domain, bukan hanya sampel.

Analogi: Untuk membuktikan bahwa ``semua angsa berwarna putih'', memeriksa jutaan angsa putih tidak cukup. Tetapi menemukan satu angsa hitam sudah cukup untuk menyangkalnya.
\end{konsep}

Kemampuan membedakan antara verifikasi empiris dan bukti formal sangat penting dalam ilmu komputer. Pengujian perangkat lunak (\textit{testing}) menunjukkan keberadaan bug (mirip kontra-contoh), tetapi test yang lulus seluruhnya tidak membuktikan ketiadaan bug --- sama seperti contoh-contoh yang mendukung tidak membuktikan pernyataan universal \cite{velleman2019}.

\section{Strategi Pemilihan Metode Bukti}

Salah satu tantangan dalam pembuktian matematis bukanlah melaksanakan langkah-langkah bukti itu sendiri, melainkan memutuskan \emph{metode mana} yang sesuai untuk digunakan. Pemilihan metode yang tepat dapat membedakan bukti yang ringkas dari rangkaian argumen yang panjang atau tidak produktif \cite{velleman2019}.

\subsection{Matriks Perbandingan Strategi}

Tabel~\ref{tab:perbandingan-strategi-bukti} menyajikan perbandingan komprehensif kelima metode pembuktian yang telah dibahas.

\begin{table}[H]
\centering
\footnotesize
\renewcommand{\arraystretch}{1.4}
\begin{tabular}{|>{\raggedright\arraybackslash}p{2.25cm}|>{\raggedright\arraybackslash}p{1.9cm}|>{\raggedright\arraybackslash}p{1.7cm}|>{\raggedright\arraybackslash}p{5.1cm}|}
\hline
\textbf{Metode} & \textbf{Asumsi Awal} & \textbf{Tujuan} & \textbf{Kapan Digunakan?} \\ \hline
\textbf{Langsung} & $p$ benar & Tunjukkan $q$ benar & Rantai deduksi dari hipotesis ke konklusi dapat disusun secara langsung \\ \hline
\textbf{Kontraposisi} & $\neg q$ benar & Tunjukkan $\neg p$ benar & Konklusi mengandung negasi; hipotesis kompleks; $\neg q$ lebih mudah diolah \\ \hline
\textbf{Kontradiksi} & $p \wedge \neg q$ benar & Temukan kontradiksi ($\bot$) & Pernyataan eksistensial negatif; irasionalitas; ketidakberhinggaan \\ \hline
\textbf{Kasus} & Partisi domain & Buktikan per kasus & Domain terbagi alami (genap/ganjil, positif/negatif/nol) \\ \hline
\textbf{Kontra-contoh} & --- & Temukan $x_0$ dengan $\neg P(x_0)$ & Menyangkal pernyataan universal \\ \hline
\end{tabular}
\caption{Matriks Perbandingan Strategi Pembuktian}
\label{tab:perbandingan-strategi-bukti}
\end{table}

\subsection{Panduan Praktis Pemilihan Metode}

Berikut panduan langkah demi langkah untuk memilih metode bukti yang tepat:

\begin{enumerate}
  \item \textbf{Identifikasi struktur pernyataan}: Apakah pernyataan merupakan implikasi ($p \rightarrow q$), biimplikasi ($p \leftrightarrow q$), atau pernyataan umum?

  \item \textbf{Coba bukti langsung terlebih dahulu}: Mulai dengan mengasumsikan hipotesis dan mencoba memanipulasi menuju konklusi. Jika langkah-langkah deduksi berjalan runtut, lanjutkan.

  \item \textbf{Jika terhambat, pertimbangkan kontraposisi}: Apakah memulai dari $\neg q$ memberikan ekspresi yang lebih mudah diolah? Apakah konklusi $q$ mengandung negasi yang menyulitkan?

  \item \textbf{Jika kedua metode di atas gagal, coba kontradiksi}: Asumsikan $p \wedge \neg q$ dan cari apakah muncul pernyataan yang mustahil. Metode ini terutama berguna untuk pernyataan tentang irasionalitas atau ketidakberhinggaan.

  \item \textbf{Pertimbangkan pembuktian kasus}: Apakah domain terbagi secara alami menjadi beberapa kategori? Apakah pernyataan lebih mudah dibuktikan per kasus?

  \item \textbf{Jika curiga pernyataan salah}: Cari kontra-contoh sebelum berusaha membuktikan.
\end{enumerate}

\subsection{Contoh Pemilihan Strategi}

\begin{contoh}
\textbf{Pernyataan}: ``Jika $n^3 + 5$ genap, maka $n$ ganjil.''

\textbf{Analisis strategi:}
\begin{itemize}
    \item \textbf{Bukti langsung?} Dari $n^3 + 5 = 2k$, didapat $n^3 = 2k - 5$. Sulit menentukan paritas $n$ dari $n^3$.
    \item \textbf{Kontraposisi?} Buktikan: ``Jika $n$ genap, maka $n^3 + 5$ ganjil.'' Jika $n = 2m$, maka $n^3 + 5 = 8m^3 + 5 = 2(4m^3 + 2) + 1$. Ini berbentuk ganjil. \checkmark
\end{itemize}

Kontraposisi memberikan bukti yang langsung dan sederhana, sedangkan pendekatan langsung menghadapi hambatan aljabar. $\blacksquare$
\end{contoh}

\subsection{Bukti Biimplikasi}

Untuk membuktikan biimplikasi $p \leftrightarrow q$ (``$p$ jika dan hanya jika $q$''), perlu dibuktikan \emph{dua} implikasi:
\begin{enumerate}
    \item $p \rightarrow q$ (arah ``jika $p$ maka $q$'')
    \item $q \rightarrow p$ (arah ``jika $q$ maka $p$'')
\end{enumerate}

Masing-masing arah dapat menggunakan metode pembuktian yang berbeda. Misalnya, arah pertama mungkin menggunakan bukti langsung, sementara arah kedua menggunakan kontraposisi. Fleksibilitas ini memungkinkan pemilihan metode yang sesuai untuk setiap arah secara independen \cite{hammack2018}.

\subsection{Hubungan dengan bab berikutnya}

Bab~13 membahas \logic{induksi matematika}, yaitu teknik untuk pernyataan pada bilangan bulat berurutan yang menggabungkan argumen pada kasus dasar dengan langkah induksi (berupa bukti langsung pada bentuk rekursif).


% ============================================================
% AKTIVITAS PEMBELAJARAN
% ============================================================

\begin{aktivitas}
  \item \textbf{Bukti langsung:} Buktikan: ``Jika $n$ bilangan genap, maka $n^2$ genap.'' Tulis setiap langkah beserta justifikasi.
  \item \textbf{Kontraposisi:} Buktikan: ``Jika $3n + 2$ genap, maka $n$ genap'' dengan kontraposisi. Jelaskan mengapa kontraposisi lebih mudah daripada bukti langsung untuk kasus ini.
  \item \textbf{Kontradiksi:} Buktikan: ``$\sqrt{3}$ irasional.'' (Petunjuk: ikuti pola bukti irasionalitas $\sqrt{2}$.)
  \item \textbf{Bukti kasus:} Buktikan: ``Untuk setiap bilangan bulat $n$, $|n| \geq 0$.'' Bagi menjadi kasus $n \geq 0$ dan $n < 0$.
  \item \textbf{Kontra-contoh:} Berikan kontra-contoh untuk menyangkal: ``Untuk setiap bilangan bulat $n > 1$, $2^n + 1$ prima.''
  \item \textbf{Pemilihan metode:} Untuk setiap pernyataan berikut, pilih metode bukti yang paling tepat, lalu buktikan atau sangkal: (a) ``Jumlah dua bilangan ganjil adalah genap.'' (b) ``Jika $n^2$ habis dibagi 3, maka $n$ habis dibagi 3.'' (c) ``Setiap bilangan yang habis dibagi 6 juga habis dibagi 12.''
\end{aktivitas}

% ============================================================
% LATIHAN DAN REFLEKSI
% ============================================================

\begin{latihan}
  \item Buktikan secara langsung: ``Jika $m$ dan $n$ bilangan ganjil, maka $m \cdot n$ ganjil.''
  \item Buktikan: ``Jika $n^2$ genap, maka $n$ genap'' dengan \textit{proof by contradiction}. Bandingkan dengan bukti kontraposisi yang dibahas di materi.
  \item Buktikan secara kontraposisi: ``Jika $n^2$ ganjil, maka $n$ ganjil.''
  \item Buktikan dengan kasus: ``Untuk setiap bilangan bulat $n$, $n^3 - n$ habis dibagi 6.'' (Petunjuk: bagi kasus berdasarkan sisa pembagian $n$ oleh 3.)
  \item Berikan kontra-contoh untuk: ``Untuk setiap bilangan real $a$ dan $b$, jika $a^2 = b^2$ maka $a = b$.''
  \item Buktikan biimplikasi: ``Bilangan bulat $n$ ganjil jika dan hanya jika $n^2$ ganjil.'' (Buktikan kedua arah.)
  \item Tentukan benar atau salah: ``Untuk setiap bilangan bulat positif $n$, $n^2 + n + 41$ prima.'' Jika benar, buktikan; jika salah, berikan kontra-contoh.
  \item \textbf{Refleksi:} (a) Kapan sebaiknya beralih dari bukti langsung ke kontraposisi? Berikan contoh konkret. (b) Dalam pengujian perangkat lunak, bagaimana kontra-contoh berkaitan dengan pencarian bug? Apakah test case yang lulus semua tes membuktikan ketiadaan bug?
\end{latihan}

% ============================================================
% ASESMEN
% ============================================================

\begin{asesmen}
\textbf{Instrumen Evaluasi Kinerja (Sub-CPMK 4.1)}

\textbf{A. Pilihan Ganda}
\begin{enumerate}
  \item Dalam pembuktian kontraposisi untuk ``Jika $P$ maka $Q$'', langkah pertama adalah:
  \begin{enumerate}
    \item Asumsikan $P$ salah
    \item Asumsikan $P$ benar
    \item Asumsikan $\neg Q$ benar
    \item Asumsikan $Q$ benar
  \end{enumerate}
  Jawaban: (c)

  \item Untuk menyangkal pernyataan ``$\forall x\, P(x)$'', cukup menunjukkan:
  \begin{enumerate}
    \item $P(x)$ salah untuk semua $x$
    \item $P(x)$ benar untuk beberapa $x$
    \item Satu kontra-contoh: $P(c)$ salah untuk suatu $c$
    \item $P(x)$ salah untuk setidaknya dua $x$
  \end{enumerate}
  Jawaban: (c)

  \item Pembuktian bahwa $\sqrt{2}$ irasional menggunakan metode:
  \begin{enumerate}
    \item Bukti langsung
    \item Bukti kontraposisi
    \item Bukti kontradiksi
    \item Bukti exhaustif
  \end{enumerate}
  Jawaban: (c)

  \item Pernyataan yang diterima benar tanpa pembuktian disebut:
  \begin{enumerate}
    \item Teorema
    \item Lemma
    \item Akibat (corollary)
    \item Aksioma
  \end{enumerate}
  Jawaban: (d)

  \item Ekuivalensi logis yang mendasari bukti kontraposisi adalah:
  \begin{enumerate}
    \item $p \rightarrow q \equiv q \rightarrow p$
    \item $p \rightarrow q \equiv \neg q \rightarrow \neg p$
    \item $p \rightarrow q \equiv \neg p \rightarrow \neg q$
    \item $p \rightarrow q \equiv p \wedge \neg q$
  \end{enumerate}
  Jawaban: (b)
\end{enumerate}

\textbf{B. Tugas Praktik}
\begin{enumerate}
  \item Buktikan: ``Jika $a$ dan $b$ bilangan rasional dengan $b \neq 0$, maka $a/b$ rasional.'' Pilih metode bukti yang tepat dan jelaskan pilihan Anda.
  \item Buktikan atau sangkal: ``Jumlah bilangan rasional dan bilangan irasional selalu irasional.'' (Petunjuk: gunakan kontradiksi.)
  \item Jelaskan perbedaan fundamental antara bukti kontraposisi dan bukti kontradiksi. Berikan contoh pernyataan yang lebih tepat dibuktikan dengan kontradiksi daripada kontraposisi.
\end{enumerate}

\textbf{Rubrik}: Merujuk Lampiran Buku.
\end{asesmen}

% ============================================================
% CHECKLIST KOMPETENSI
% ============================================================

\begin{checklist}
  \item Saya memahami pengertian teorema, bukti, aksioma, lemma, akibat, dan konjektur.
  \item Saya mampu menyusun bukti langsung untuk implikasi $p \rightarrow q$ dengan justifikasi pada setiap langkah.
  \item Saya memahami ekuivalensi $p \rightarrow q \equiv \neg q \rightarrow \neg p$ dan mampu menerapkan bukti kontraposisi.
  \item Saya mampu menyusun bukti kontradiksi dengan mengasumsikan negasi dan menurunkan kemustahilan.
  \item Saya memahami dan mampu menerapkan pembuktian dengan kasus serta pembuktian exhaustif.
  \item Saya mampu menyangkal pernyataan universal dengan menemukan kontra-contoh yang valid.
  \item Saya mampu memilih metode bukti yang tepat berdasarkan struktur pernyataan yang hendak dibuktikan.
  \item Saya memahami perbedaan antara verifikasi empiris dan pembuktian formal.
  \item Saya mampu membuktikan biimplikasi ($p \leftrightarrow q$) dengan membuktikan kedua arah implikasi.
\end{checklist}

% ============================================================
% RANGKUMAN
% ============================================================

\begin{rangkuman}
Bab ini membahas berbagai metode pembuktian matematis yang menjadi landasan bagi penalaran formal dalam matematika dan ilmu komputer.

\textbf{Poin Kunci:}
\begin{itemize}
  \item \textbf{Teorema} adalah pernyataan yang telah dibuktikan, \textbf{bukti} adalah demonstrasi deduktif kebenarannya. Terminologi terkait meliputi aksioma (kebenaran tanpa bukti), lemma (teorema bantu), akibat (konsekuensi langsung), dan konjektur (dugaan belum terbukti).
  \item \textbf{Bukti langsung} mengasumsikan $p$ benar dan menurunkan $q$ benar melalui manipulasi aljabar dan definisi. Ini adalah metode yang paling intuitif dan menjadi pilihan pertama.
  \item \textbf{Bukti kontraposisi} membuktikan $\neg q \rightarrow \neg p$ yang ekuivalen dengan $p \rightarrow q$. Berguna saat konklusi mengandung negasi atau hipotesis terlalu kompleks untuk diolah langsung.
  \item \textbf{Bukti kontradiksi} mengasumsikan negasi dari pernyataan dan menurunkan kemustahilan. Metode paling kuat dan serbaguna, terutama untuk pernyataan tentang irasionalitas dan ketidakberhinggaan.
  \item \textbf{Bukti kasus} mempartisi domain menjadi kasus-kasus yang saling lepas dan lengkap, lalu membuktikan per kasus. Pembuktian exhaustif memeriksa setiap elemen domain berhingga.
  \item \textbf{Kontra-contoh} menyangkal pernyataan universal dengan menunjukkan satu elemen yang tidak memenuhi. Satu kontra-contoh cukup, tetapi ribuan contoh pendukung tidak membuktikan kebenaran universal.
\end{itemize}

\textbf{Kata Kunci}: teorema, bukti, aksioma, lemma, akibat, konjektur, bukti langsung, kontraposisi, kontradiksi, bukti kasus, pembuktian exhaustif, kontra-contoh, biimplikasi
\end{rangkuman}

\ifSubfilesClassLoaded{
  \renewcommand{\bibname}{Daftar Pustaka}
  \bibliographystyle{plain}
  \bibliography{../references}
}{}
\end{document}
